\chapter{Rejestr stanu książki}
\markboth{Rejestr stanu książki}{Rejestr stanu książki}
Liczby w tym dodatku są \emph{liczone} z danych i źródeł (\code{tools/gen.py}), nie wpisywane.

\begin{center}\small
\begin{tabular}{@{}lr@{}}
\toprule
rozdziały w planie & \val{chapters.total}\\
rozdziały napisane & \val{written.total}\\
ćwiczenia z kluczem & \val{exercises.total}\\
zadania na tydzień & \val{tasks.total}\\
ramki \emph{Pułapka} & \val{traps.total}\\
twierdzenia oznaczone \code{\textbackslash verify} (wymagają źródła) & \val{verify.total}\\
\bottomrule
\end{tabular}
\end{center}

\section*{Dług, który książka ujawnia}
\begin{itemize}
\item Recenzja lekarza medycyny snu: NIEUSTALONA
\item Liczby oznaczone \verify{tak} (pomarańczowe [?]) to twierdzenia, których autor nie sprawdził w pierwotnym źródle w tej sesji.
\item Liczby bez oznaczenia są albo policzone z danych w repozytorium, albo są arytmetyką pokazaną w tekście.
\end{itemize}
